\section{Functional CALM}
\label{sec:functional-calm}

The CALM Theorem was proved for queries in a logic programming language,
evaluated by relational transducers, and its applicability outside that
setting was never made precise.  Here we use a simpler model that applies to
any specification expressible as an input/output relation over ordered domains.  A
specification describes allowed input/output
behavior without committing to a language, placement, protocol, or message
schedule.  This section addresses the case that prior work on both CALM and
CRDTs shared: specifications that are functions, in which each
input has a single specified output.  Section~\ref{sec:election-story} walks
through the framework informally on an election.
Section~\ref{sec:functional-specifications} defines functional specifications
and monotonicity, Section~\ref{sec:determination} defines the transition system
and Consistency, Section~\ref{sec:functional-calm-theorem} proves Functional
CALM, and Section~\ref{sec:canonical-transition-system} identifies when an
output can be emitted without commitment.

\subsection{The framework on an election}
\label{sec:election-story}

We begin with an informal, example-driven treatment.
Terms introduced in italics are defined formally in the rest of this section
and in Sections~\ref{sec:distributed-execution}--\ref{sec:complete-calm};
italics also mark the occasional
standard term borrowed from database theory.  Nothing here is needed for the
proofs, but every notion the Functional CALM theorem rests on---specification
states, input events, retirement, commitments, upper fibers, determination,
output actions, runs, and consistency---is exercised on the example.

The framework begins with an \emph{ordered specification}
$\Spec=(I,\InputOrd,O,\OutputOrd,S)$: a table $S\subseteq I\times O$ that
captures exactly which outputs are within spec at each input, with a partial
order on each column.  In this section, input functionally determines output:
every input is paired with exactly one output.
Section~\ref{sec:ambiguous-specifications} later drops this functional
dependency.

\begin{figure*}[t]
  \centering
  \small
  \newcommand{\rowhdr}{\toprule Input & Output \\ \midrule}
  \setlength{\tabcolsep}{6pt}
  \begin{tabular}[t]{@{}l c@{}}
    \multicolumn{2}{@{}l}{(a) Initial state, current input $x\cdot\cdot$}\\
    \rowhdr
    \sh{$\mathbf{x\cdot\cdot}$} & \sh{$\bot$} \\
    $xy\cdot$ & \sh{$\bot$} \\
    $xyx$ & \sh{$x$} \\
    $xyy$ & \sh{$y$} \\
    $xx\cdot$ & \sh{$x$} \\
    $xxx$ & \sh{$x$} \\
    $xxy$ & \sh{$x$} \\
    \bottomrule
    \multicolumn{2}{@{}l}{\footnotesize possible outputs $\{\bot,x,y\}$; undetermined}\\
    \multicolumn{2}{@{}l}{\footnotesize $\bot$ needs no commitment}
  \end{tabular}
  \hfill
  \begin{tabular}[t]{@{}l c@{}}
    \multicolumn{2}{@{}l}{(b) After $\mathsf{input}(xy\cdot)$}\\
    \rowhdr
    \retired{$x\cdot\cdot$} & \retired{$\bot$} \\
    \sh{$\mathbf{xy\cdot}$} & \sh{$\bot$} \\
    $xyx$ & \sh{$x$} \\
    $xyy$ & \sh{$y$} \\
    \retired{$xx\cdot$} & \retired{$x$} \\
    \retired{$xxx$} & \retired{$x$} \\
    \retired{$xxy$} & \retired{$x$} \\
    \bottomrule
    \multicolumn{2}{@{}l}{\footnotesize possible outputs $\{\bot,x,y\}$; undetermined}\\
    \multicolumn{2}{@{}l}{\footnotesize $\bot$ needs no commitment}
  \end{tabular}
  \hfill
  \begin{tabular}[t]{@{}l c@{}}
    \multicolumn{2}{@{}l}{(c) After $\mathsf{input}(xx\cdot)$ instead}\\
    \rowhdr
    \retired{$x\cdot\cdot$} & \retired{$\bot$} \\
    \retired{$xy\cdot$} & \retired{$\bot$} \\
    \retired{$xyx$} & \retired{$x$} \\
    \retired{$xyy$} & \retired{$y$} \\
    \sh{$\mathbf{xx\cdot}$} & \sh{$x$} \\
    $xxx$ & \sh{$x$} \\
    $xxy$ & \sh{$x$} \\
    \bottomrule
    \multicolumn{2}{@{}l}{\footnotesize possible outputs $\{x\}$; determined,}\\
    \multicolumn{2}{@{}l}{\footnotesize input-resolved; $x$ needs no commitment}
  \end{tabular}
  \hfill
  \begin{tabular}[t]{@{}l c@{}}
    \multicolumn{2}{@{}l}{(d) (b), then polls closed}\\
    \rowhdr
    \retired{$x\cdot\cdot$} & \retired{$\bot$} \\
    \sh{$\mathbf{xy\cdot}$} & \sh{$\bot$} \\
    \retired{$xyx$} & \retired{$x$} \\ \strikerow
    \retired{$xyy$} & \retired{$y$} \\ \strikerow
    \retired{$xx\cdot$} & \retired{$x$} \\
    \retired{$xxx$} & \retired{$x$} \\
    \retired{$xxy$} & \retired{$x$} \\
    \bottomrule
    \multicolumn{2}{@{}l}{\footnotesize possible outputs $\{\bot\}$; sealed and}\\
    \multicolumn{2}{@{}l}{\footnotesize determined by commitment; $\bot$ is the result}
  \end{tabular}
  \caption{The election specification and three states reached from it.  Only rows above the input $x\cdot\cdot$ are shown.  The current input is
  in bold; shaded cells are the current input and outputs still possible at it.
  Gray: the row has been retired by an input event.
  Struck through: the row has been removed by a commitment; no remaining row
  preserves its contents.  Panels (b) and (c) follow input events; panel (d)
  follows an output action for $\bot$ that contains a seal commitment.}
  \label{fig:election-frames}
\end{figure*}

\paragraph{A specification is a table.}
A specification table lists the output paired with each input.  It need not be finite. Squaring, for example,
has rows $(2,4)$, $(3,9)$, $(4,16),\ldots$.  Everything that follows is a
matter of removing rows from an initial specification.

Figure~\ref{fig:election-frames}(a) shows a set of rows taken from the specification of a small election.
Three voters $v_1,v_2,v_3$ each cast at most one immutable ballot for $x$ or
$y$; the result is whichever candidate has two ballots, and $\bot$ (no result
yet) until one does.  An input represents the
ballots seen so far, written as a three-letter string whose $k$th position holds
$v_k$'s ballot or a dot if it has not arrived; only rows above $x\cdot\cdot$,
where $v_1$ has voted $x$, are shown.  The first row says that with $v_1$'s
ballot alone the result is $\bot$; row $xx\cdot$ says that once $v_2$ also votes
$x$ the result is $x$; row $xyy$ says that if $v_2$ and $v_3$ both vote $y$ the
result is $y$.

The rows are only part of the specification.  Each column also carries a
partial order, and both orders are design decisions.  The \emph{input order}
says what counts as more input: here, more ballots, so $xy\cdot$ \emph{refines}
$x\cdot\cdot$, while $xx\cdot$ and $xy\cdot$ are incomparable.  The \emph{output
order} says when a later output keeps every promise an earlier one made: here
$\bot$ lies below both $x$ and $y$, so ``no result yet'' can eventually refine
to a winner, but $x$ and $y$ are incomparable, so neither winner refines the
other.  (A
different design could make all outputs incomparable, so that every result is
final; we return to that choice below.)  Given these orders, one can follow any input
upward through its refinements, say $x\cdot\cdot$ to $xy\cdot$ to $xyx$: the
output climbs too, from $\bot$ to $x$, and never steps down or sideways.  That
is what it means for the \emph{specification} to be \emph{monotone}, a property
of the rows together with the orders.  Viewing the specification as a function
$f:I\to O$, monotonicity says that
$i\InputOrd j\Longrightarrow f(i)\OutputOrd f(j)$.

\paragraph{Input arrives.}
Now $v_2$'s ballot arrives, for $y$; panel (b) shows the result.  The
\emph{current input}, the ballots seen so far, has advanced from $x\cdot\cdot$
to $xy\cdot$, shown in bold.  One rule governs what an input event does to the specification:
only the rows at the current input and its refinements remain; every other row is
\emph{retired}---removed from the resulting remainder table.  Each state pairs the current input
with the rows that remain: it records what
is still possible and how far input has come.  Each panel shows one such state,
with its remainder table and current input in bold.  The inputs that still have rows are the \emph{valid} ones.
The environment supplies new input; the system passively retires rows that are no
longer reachable.  Here $v_2$'s ballot
extends the cumulative input from $x\cdot\cdot$ to $xy\cdot$, and the system
records that change as $\mathsf{input}(xy\cdot)$.
The state depends only on the input reached, not on the order or
grouping in which the ballots arrived.  In panel (b) the input has grown
from $x\cdot\cdot$ to $xy\cdot$, so the row at $x\cdot\cdot$ is retired; the
rows at $xx\cdot$, $xxx$, and $xxy$ are retired because no input refines both
$xx\cdot$ and $xy\cdot$.  The retired output $\bot$ at $x\cdot\cdot$ is
refined by the $\bot$ at $xy\cdot$, so growing the input lost nothing the
smaller input's output claimed; this is a consequence of the monotonicity
observed above.

\paragraph{What could still happen.}
While input may still grow, what matters is not only the current row but what
any input that could still arrive might require.  The future of the election
is its current input together with its valid input refinements, and the outputs of
those rows form the current input's \emph{upper fiber}.  In
panel (b) those rows are the current row and the two refinements $xyx$ and
$xyy$ in which $v_3$ votes, so the upper fiber, shaded, is $\{\bot,x,y\}$:
either candidate could still win.

Panel (c) shows the alternative in which $v_2$ voted $x$.  The rows at or
above $xx\cdot$ all say $x$, so the upper fiber is the singleton $\{x\}$ and
the current state is \emph{determined}: the arriving input resolves the
remaining ambiguity by ruling out ballot histories that do not extend the
ballots received.  No commitment was needed; the result is input-resolved.
Below we return to panel
(b), where several possible outputs remain and the system must commit to reduce
them.

\paragraph{Commitments.}
Return to panel (b).  Consider a registrar who closes the polls before
exposing $\bot$, so any later ballot from $v_3$ will not count.  Panel (d)
removes the rows at $xyx$ and $xyy$ while
leaving the current input $xy\cdot$ unchanged.  The system has thereby made
those input refinements invalid.  This irrevocable act is a \emph{commitment},
specifically a \emph{seal commitment}: the system, rather than new input from
the environment, has ruled out behavior that was within spec.

The specification does not define a notion of commitments; that authority belongs to a second
component of our model: the system's \emph{commitment basis} $\mathcal B$, which defines
how the system may irrevocably remove rows.  A system whose polls can close includes this commitment in
its basis; one whose polls cannot close does not.  After this output action, no refinement of $xy\cdot$ remains valid, its upper
fiber is $\{\bot\}$, and the commitments have determined the result.

Because a functional specification has one row per input, any commitment that
removes a row makes that input invalid.  Thus every commitment in this section
is a seal commitment; it may invalidate one future input or many.

\begin{figure}[t]
  \centering
  \small
  \setlength{\tabcolsep}{6pt}
  \begin{tabular}[t]{@{}l c@{}}
    \multicolumn{2}{@{}l}{(e) As (a), all outputs incomparable, polls closed}\\
    \toprule
    Input & Output \\ \midrule
    \sh{$\mathbf{x\cdot\cdot}$} & \sh{$\bot$} \\
    \retired{$xy\cdot$} & \retired{$\bot$} \\ \strikerow
    \retired{$xyx$} & \retired{$x$} \\ \strikerow
    \retired{$xyy$} & \retired{$y$} \\ \strikerow
    \retired{$xx\cdot$} & \retired{$x$} \\ \strikerow
    \retired{$xxx$} & \retired{$x$} \\ \strikerow
    \retired{$xxy$} & \retired{$x$} \\ \strikerow
    \bottomrule
    \multicolumn{2}{@{}l}{\footnotesize possible outputs $\{\bot\}$; determined}\\
    \multicolumn{2}{@{}l}{\footnotesize by commitment; $\bot$ needs no further commitment}
  \end{tabular}
  \caption{Panel (a) of Figure~\ref{fig:election-frames} with all outputs
  incomparable, after the registrar emits $\bot$ while closing the polls at
  $x\cdot\cdot$.  Every refinement of $x\cdot\cdot$ is removed by the output
  action's seal commitment.  Legend as in
  Figure~\ref{fig:election-frames}.}
  \label{fig:nonmonotone-frame}
\end{figure}

\paragraph{Output actions.}
In a functional specification, $o$ is \emph{exposable} when every output still
allowed at the current input or one of its valid refinements preserves $o$.  If
this test fails, emitting $o$ requires an authorized commitment that removes
the conflicting rows.  In Figure~\ref{fig:election-frames}, $\bot$ is
exposable in panels (a), (b), and (d), while $x$ is exposable in panel (c).

The transition system records one atomic $\mathsf{output}(o)$ action.  The
action prepares $o$ for exposure by issuing any required
commitments and retiring rows that $o$ supersedes, then emits $o$.
The internal stages of this atomic action are not separate transitions: no event
can intervene.  The output emission itself remains explicit
because the environment may use the emitted result to choose later input.

\paragraph{When monotonicity fails.}
Now change one thing in our example: make all outputs incomparable, dropping
$\bot\OutputOrd x,y$, so that an output refines only itself; this is the
\emph{equality order}.  The rows are the same, but in panels (a) and (b)
nothing is exposable: $\bot$ no longer lies below the $x$ and $y$ in the upper
fiber.  Emitting any output requires a seal commitment.  Figure~\ref{fig:nonmonotone-frame} shows the result of the output
action for $\bot$ at $x\cdot\cdot$: its commitment removes every refinement,
and the upper fiber collapses to $\{\bot\}$.  That contrast is Functional
CALM (Section~\ref{sec:functional-calm}): a monotone functional specification
can emit its outputs without commitment; a non-monotone one requires a
commitment in some output action.

\paragraph{A transition system.}
We now have the ingredients for a state transition model.  A state consists of
two things: the current input and the rows of the specification that remain.
Each panel above depicts one such state, and each action moves from one state to
another.  An input event records a larger cumulative input and retires rows
that it has passed.  An output action keeps the input fixed, performs any
required commitment, emits its result, and retires rows superseded by that
result.  Inputs cross the interface from the environment; outputs cross it from
the system.  Panels (b) through (e) are reached from panel (a) by input or
output actions.  The formal development uses this one transition system
throughout.

\paragraph{Consistency.}
Our goal is \emph{confluence}: any finite collection of input/output paths
through the transition system can be brought to one common state, whatever
further input arrives.  We call a specification \emph{consistent} when all
such paths are confluent.

Our framework can encode either a general consistency model that spans many
applications or an application-specific requirement such as an invariant.  In
the election, the table and its orders license the application's legal
behaviors: which results it permits, and which later results preserve earlier
ones.  Consistency then asks whether paths through those behaviors can still
meet.  Appendix~\ref{app:registers} applies the framework to linearizability
and causal consistency.  These application-independent consistency models
license different register behaviors and therefore yield different confluence
results.  Section~\ref{sec:consistency-as-confluence} develops this distinction
generally.  Specifications can also be composed: joining component
tables on their shared interface yields an end-to-end specification
(Appendix~\ref{app:composition}).

Monotonicity gives a condition directly on the specification table: whenever
input grows from $i$ to $k$, the output at $k$ refines the output at $i$.
Functional CALM proves that this condition is equivalent to confluence: a
functional specification is consistent exactly when it is monotone.

The rest of this section states these notions formally.
Figure~\ref{fig:formalism-summary} collects the formal objects of Functional
CALM; Figure~\ref{fig:complete-formalism-summary} will collect the additions
for Complete CALM.

\begin{figure*}[t]
  \centering
  \small
  \setlength{\tabcolsep}{6pt}
  \renewcommand{\arraystretch}{1.25}
  \begin{tabular}{@{}l l l@{}}
    \toprule
    \textbf{Specification} & & \\
    \midrule
    specification & $f:(I,\InputOrd)\to(O,\OutputOrd)$; \quad $S=\{(i,f(i)):i\in I\}$ & Def.~\ref{def:functional-specification} \\
    monotone (functional) & $i\InputOrd j\ \Longrightarrow\ f(i)\OutputOrd f(j)$ & Def.~\ref{def:functional-monotonicity} \\
    \addlinespace
    \textbf{Shared state and input} & & \\
    \midrule
    specification state & $(i,R)$: \ current input $i$ and remainder table $R\subseteq S$ & Def.~\ref{def:residual-state} \\
    initial state & $\sigma_0(i)=(i,\{(k,o)\in S:i\InputOrd k\})$ & Eq.~\ref{eq:canonical-initial-state} \\
    input event & $(i,R)\xrightarrow{\mathsf{input}(j)}(j,R\mathbin{\upharpoonright_I}j)$ & Def.~\ref{def:io-transition-system} \\
    \addlinespace
    \textbf{Output handling} & & \\
    \midrule
    output action & commitments and retirements, then emit $o$ & Eq.~\ref{eq:output-action} \\
    canonical commitment basis & supports every permitted output action & Sec.~\ref{sec:determination} \\
    exposable output & output action uses no commitment & Def.~\ref{def:exposable} \\
    run & finite input and output actions & Def.~\ref{def:run} \\
    \addlinespace
    \textbf{Results} & & \\
    \midrule
    consistent & paths confluent: \ $\bigcap_r R_r(j)\neq\varnothing$ at every common refinement $j$ of their current inputs & Def.~\ref{def:consistency}, Prop.~\ref{prop:confluence} \\
    Functional CALM & consistent $\iff$ monotone & Thm.~\ref{thm:functional-calm} \\
    \bottomrule
  \end{tabular}
  \caption{Reference for Functional CALM: one transition system embeds the
  declarative specification and commitment basis in atomic input/output actions.}
  \label{fig:formalism-summary}
\end{figure*}

\subsection{Functional specifications and monotonicity}
\label{sec:functional-specifications}

\begin{definition}[Functional ordered specification]
\label{def:functional-specification}
A \emph{functional ordered specification} is a function
        \begin{equation}
  f:(I,\InputOrd)\longrightarrow(O,\OutputOrd),
  \label{eq:functional-specification}
        \end{equation}
where $(I,\InputOrd)$ and $(O,\OutputOrd)$ are partially ordered sets.  We call
$\InputOrd$ the \emph{input-refinement order}: $i\InputOrd j$ means that $j$
contains at least the input information in $i$.  We call $\OutputOrd$ the
\emph{output-refinement order}: $o\OutputOrd u$ means that $u$ preserves every
externally visible claim in $o$ and may add more.
\end{definition}

These are two separate partial orders.  The input order compares inputs; the
output order compares outputs.  Neither orders executions or rows in a table.
Either may be the \emph{equality order}, under which an element refines only
itself; on outputs this is the case in which no output can be extended, only
replaced.

\begin{definition}[Monotonicity]
\label{def:functional-monotonicity}
A functional ordered specification $f$ is \emph{monotone} if
        \begin{equation}
  i\InputOrd j
  \quad\Longrightarrow\quad
  f(i)\OutputOrd f(j).
  \label{eq:functional-monotonicity}
        \end{equation}
\end{definition}

This is the standard mathematical definition of a monotone function.  In
systems terms, an output exposed from partial input remains valid as more input
arrives.

\subsection{The transition system}
\label{sec:determination}

Section~\ref{sec:election-story} previewed the transition model that embeds
the declarative specification and a system commitment basis in a simple I/O
automaton.  Its state is a specification state.  An input action records a
cumulative input supplied by the environment.  An atomic output action prepares
its result by removing rows for reasons of commitment or retirement, then emits
that result.  The action records whether it used a commitment, and the
environment may use its emitted output to choose a later input.

\paragraph{The specification and remainder tables.}
The relation $S\subseteq I\times O$, with the functional dependency that
inputs determine outputs, is the \emph{specification table}; its rows
are the pairs $(i,f(i))$, one at each input.  During a run, $R\subseteq S$ is
the \emph{remainder table}; it records the rows not yet removed.  No action
adds a row.

Inputs are \emph{cumulative}: an input is everything seen so far, and
$i\InputOrd j$ means that $j$ contains at least the information in $i$.  This
is a requirement on representation, not on application-level effect.  Every
change must be represented by adding information to the input, even when its
meaning invalidates an earlier conclusion.  Thus a correction records both the
original assertion and its correction, a retraction records the assertion and
its withdrawal, and a conflict records both sides.  The framework otherwise
uses only the elements of $I$, their refinement order, and the outputs paired
with each input.  Write
$k\prec_I i$ when $k\InputOrd i$ and $k\neq i$, and
$\mathord\uparrow_I i=\{j\in I:i\InputOrd j\}$ for the refinements of $i$.

\paragraph{States.}
A state pairs the current input with a remainder table.  Every panel of
Figure~\ref{fig:election-frames} is such a state.

\begin{definition}[Specification state]
\label{def:residual-state}
A \emph{specification state} is a pair $\sigma=(i,R)$ such that
\begin{enumerate}
  \item $i\in I$ is the \emph{current input};
  \item $R\subseteq S$;
  \item $R$ has a row at the current input:
        \begin{equation}
          (i,o)\in R\ \text{for some }o;
          \label{eq:current-input-validity}
        \end{equation}
  \item only refinements of the current input have rows in $R$:
        \begin{equation}
          (k,o)\in R
          \quad\Longrightarrow\quad i\InputOrd k.
          \label{eq:retired-input-absence}
        \end{equation}
\end{enumerate}
An input $k$ is \emph{valid} when $R$ has a row at $k$, and otherwise
\emph{invalid}.
\end{definition}

Valid inputs above the current input are its still-admissible refinements.  For
an input $i$, the \emph{canonical initial state} is
        \begin{equation}
  \sigma_0(i)=\bigl(i,\ \{(k,o)\in S:\ i\InputOrd k\}\bigr),
  \label{eq:canonical-initial-state}
        \end{equation}
the specification table restricted to $i$ and its refinements.
Figure~\ref{fig:election-frames}(a) shows $\sigma_0(x\cdot\cdot)$.

\paragraph{Input events.}
For $j\in I$ with $i\InputOrd j$, an input event makes $j$ current and removes
every row whose input does not refine $j$:
        \begin{equation}
  (i,R)\xrightarrow{\mathsf{input}(j)}
  \bigl(j,\ R\mathbin{\upharpoonright_I}j\bigr),
  \qquad
  R\mathbin{\upharpoonright_I}j=\{(k,u)\in R:\ j\InputOrd k\},
  \label{eq:input-event}
        \end{equation}
defined when $R$ has a row at $j$.  The event retires every row whose input
does not refine the new current input $j$.

The event label is already cumulative: it supplies the new representation $j$,
not a raw fragment that this transition system must fold into $i$.  How an
environment or implementation computes $j$ lies outside the specification.  It
may use union, a semilattice join, or another rule.  If two pieces of input have
no common refinement, this input domain cannot represent their joint receipt;
the interface must exclude that combination or use a richer cumulative
representation.  If several common refinements exist, the surrounding system
may choose one and supply it as the next input.  The framework requires neither
a least common refinement nor a lattice.

Input events compose once their final target is fixed:
$\mathsf{input}(j)$ followed by $\mathsf{input}(k)$ has the same effect as
$\mathsf{input}(k)$ alone, and $\mathsf{input}(i)$ at $(i,R)$ is the identity.

\paragraph{Commitments and output handling.}
For $(i,o)\in R$, the final table needed to honor emission of $o$ is
\begin{equation}
  R_o=R\mathbin{\upharpoonright_O}o
      =\{(k,u)\in R:o\OutputOrd u\}.
  \label{eq:residual-refining-fibers}
\end{equation}
In a functional specification each input has one row.  Consequently every row
in $R\setminus R_o$ lies at a valid refinement whose only output conflicts with
$o$.  Emitting $o$ must make each such input invalid.

A \emph{commitment basis} $\mathcal B$ is a set of partial functions on
specification states.  Each commitment keeps the current input fixed and
strictly decreases the remainder table.  A \emph{seal commitment} removes every
row at one or more valid noncurrent inputs, making those inputs invalid.  For a
functional specification, these are the only commitments needed to emit an
output.

For every specification state and $(i,o)\in R$, the canonical seal commitment
is the partial function
\begin{equation}
  s_{(i,R),o}(i,R)=(i,R_o),
  \label{eq:canonical-seal-commitment}
\end{equation}
when $R_o\neq R$; identity cases are omitted.  The \emph{canonical commitment
basis} $\mathcal B_S^{\mathrm{can}}$ contains every such commitment.

An output action for $o$ applies a finite, possibly empty sequence of commitments
from the system's basis, retires any remaining rows superseded by $o$, and emits
$o$; these effects are atomic.  The action itself records which commitments it
applied, although its interface label is simply $\mathsf{output}(o)$.  In the
functional case there are no output retirements, so the applied commitments
must take $(i,R)$ to $(i,R_o)$.  Under the canonical basis,
\begin{equation}
  (i,R)\xrightarrow{\mathsf{output}(o)}(i,R_o).
  \label{eq:output-action}
\end{equation}

The canonical basis supports every output permitted by the specification.  This
is essential for CALM: a weaker basis could hide non-monotonicity merely by
making a conflicting output unavailable.  Another basis supports an output action
for $o$ when its commitments reach some $(i,R')$ with
$(i,o)\in R'\subseteq R_o$; it may therefore remove more than the canonical
action.  Exact equality $R'=R_o$ is required only of the canonical action.  Section~\ref{sec:safe-exposure} extends
output actions and retirement to specifications with several outputs at one
input.

\begin{definition}[Transition system and run]
\label{def:io-transition-system}
\label{def:transition-system}
The \emph{transition system} has specification states as its states, input
events as in Equation~\ref{eq:input-event}, and atomic output actions as in
Equation~\ref{eq:output-action}.  A \emph{run} is any finite transition sequence
from a canonical initial state $\sigma_0(i)$.  Unless stated otherwise, runs use
the canonical commitment basis.
\label{def:run}
\end{definition}

A specification state is \emph{run-generated} when it is reachable via a run.
This is a provenance condition, not another local invariant on a remainder
table.  In this functional section, missing rows arise from input retirements
and seal commitments.  Section~\ref{sec:safe-exposure} revisits this account for
specifications with several outputs at one input.

\begin{lemma}[Rows remaining after a run]
\label{lem:run-residual}
A run starts at $\sigma_0(i_0)$, ends at $(i,R)$, and contains output actions
emitting $o_1,\ldots,o_m$.  Its final state is
        \begin{equation}
  \left(i,\ \left\{(k,x)\in S:
    \begin{array}{l}
      i\InputOrd k,\ \text{and}\\
      o_\ell\OutputOrd x\ \text{for every }\ell
    \end{array}
  \right\}\right).
  \label{eq:run-residual}
        \end{equation}
\end{lemma}

\begin{proof}
Each input event intersects $R$ with $\{(k,x):j\InputOrd k\}$ and each
canonical output action with $\{(k,x):o\OutputOrd x\}$.  In the functional
case, its commitments have exactly this effect and output retirement removes no
additional row.  These intersections commute, and the input
restrictions compose to the restriction at the final input $i$.
\end{proof}

A \emph{common refinement} of inputs $i_1,\ldots,i_k$ is an input $j$ with
$i_r\InputOrd j$ for every $r$.

\begin{definition}[Consistency]
\label{def:consistency}
An ordered specification is \emph{consistent} when every finite family of runs
under its canonical commitment basis is confluent: if the runs end at states
$(i_1,R_1),\ldots,(i_k,R_k)$ and $j$ is a common refinement of their current
inputs, one state with current input $j$ is reachable by input and output
actions from each endpoint.  The runs are then \emph{confluent at} $j$.
\end{definition}

Equivalently, Consistency quantifies over every finite family of run-generated
endpoints.  This covers every run-generated starting state, since any path from
one extends the run that generated it.  It does not quantify over arbitrary
specification states or arbitrary commitments.  The canonical commitment basis
includes exactly the output-directed commitments needed to exercise every
output permitted by the specification; a weaker basis could make a
non-monotone specification appear consistent by disabling its conflicting
outputs.  A coordinating component's exported interface is tested separately
as a specification in Appendix~\ref{app:composition}.

\begin{proposition}[Confluence at a common refinement]
\label{prop:confluence}
Runs ending at states $(i_1,R_1),\ldots,(i_k,R_k)$ are confluent at $j$ if and
only if some row at $j$ survives in every final state,
        \begin{equation}
  (j,u)\in R_r\quad\text{for every }r,\ \text{for some }u\in O.
  \label{eq:reconciliation}
        \end{equation}
For any such $u$, $\mathsf{input}(j)$ followed by $\mathsf{output}(u)$ takes
every endpoint to the state
$\bigl(j,\{(k,x)\in S:j\InputOrd k,\ u\OutputOrd x\}\bigr)$.
\end{proposition}

\begin{proof}
A common state with current input $j$ has a row at $j$, and that row belongs to
every $R_r$.  Conversely, let $(j,u)$ belong to every $R_r$.  Then
$\mathsf{input}(j)$ is defined at every $(i_r,R_r)$.  By
Lemma~\ref{lem:run-residual}, $(j,u)\in R_r$ says that $u$ refines every output
emitted on run $r$.  After $\mathsf{input}(j)$ and $\mathsf{output}(u)$, the
same lemma leaves exactly the rows of $S$ above $j$ and $u$, independent of the
starting run.
\end{proof}

A single run can witness inconsistency: an output action for one permitted
output may remove every row at a refinement of its current input.
Consistency describes
no protocol and promises no delivery or termination;
Section~\ref{sec:discussion} returns to this.

\subsection{The Functional CALM theorem}
\label{sec:functional-calm-theorem}

\begin{theorem}[Functional CALM]
\label{thm:functional-calm}
A functional ordered specification is consistent if and only if it is
monotone.
\end{theorem}

\begin{proof}
Every input has one row.  By Lemma~\ref{lem:run-residual}, a run ending at
input $i$ retains $(k,f(k))$ at $k\succeq_I i$ exactly when every output
$f(i_\ell)$ emitted on the run satisfies $f(i_\ell)\OutputOrd f(k)$.  By
Proposition~\ref{prop:confluence}, a finite family of runs is confluent at a
common refinement $k$ exactly when every output emitted on any run lies below
$f(k)$.  If $f$ is monotone this holds because each emitted $f(i_\ell)$ has
$i_\ell\InputOrd k$.  If $f$ is not monotone, some $i\prec_I k$ has
$f(i)\not\OutputOrd f(k)$; the run that starts at $\sigma_0(i)$ and emits
$f(i)$ has no row at $k$ and cannot advance to its own refinement $k$.
\end{proof}

This is the familiar CALM dichotomy stated at the specification interface
rather than inside a particular programming language.  Classical CALM fits the
pattern with inputs $J$ ordered by fact inclusion and outputs $Q(J)$ ordered by
answer inclusion.  Functional monotonicity becomes
        \begin{equation}
  J\subseteq J'
  \quad\Longrightarrow\quad
  Q(J)\subseteq Q(J'),
  \label{eq:query-monotonicity}
        \end{equation}
the standard definition of a monotone query~\cite{abiteboul1995foundations},
which positive Datalog satisfies and negation may not.  Functional CALM
extends the statement from fact sets under inclusion to arbitrary input and
output refinement orders.  The usual CRDT state interface has the same
functional shape, with received payloads ordered by inclusion and states by
the semilattice order~\cite{shapiro2011crdt};
Sections~\ref{sec:query-instance} and~\ref{sec:crdt-instance} take up both
instances.

\subsection{Output actions and exposability}
\label{sec:canonical-transition-system}

The canonical commitment basis exists to exercise every output the specification
permits, not to claim that every implementation supports every output action.
An implementation supplies its own basis; it can take an output action only
when commitments in that basis, together with permitted retirements, produce
the required effect.  The following definitions identify when no commitment is
required.

For $R\subseteq S$, the \emph{upper fiber} at $i$ collects outputs at $i$ and
its still-valid refinements:
        \begin{equation}
  R^\uparrow(i)=\{o\in O:\exists j\in I,\ i\InputOrd j\text{ and }(j,o)\in R\}.
  \label{eq:upper-fiber}
        \end{equation}
A state is \emph{determined} when $|R^\uparrow(i)|=1$, regardless of how it was
reached.  \emph{Determination} is movement toward or arrival at that condition.
A commitment may advance determination by excluding an alternative that
otherwise remains possible; an input event may instead resolve ambiguity by
retiring histories that no longer extend the input received.  We call the
resulting outcomes \emph{commitment-resolved} and \emph{input-resolved},
respectively.

\begin{definition}[Exposable output, functional case]
\label{def:exposable}
For a functional specification, $(i,o)\in R$ is \emph{exposable} at $(i,R)$
when its output action requires no commitment.  Equivalently, the output action
removes no row: $o\OutputOrd f(j)$ at every valid $j\succeq_I i$.
\end{definition}

A state is \emph{sealed} when no strict refinement of its current input remains
valid.  A sealed functional state is determined.  Commitments that produce such
a state make it commitment-resolved, although determination need not seal every
strict refinement at once.

A system's \emph{strategy} chooses input and output actions, including which
commitments from its basis an output action applies, using the specification
and any additional strategy state.  Strategy state is not part of the
specification state and affects the remainder table only through the selected
transitions.

\begin{proposition}[Functional exposure boundary]
\label{prop:functional-exposure-boundary}
A functional ordered specification is monotone if and only if every output at
every state reachable without commitment is exposable.
\end{proposition}

\begin{proof}
If $f$ is monotone, a commitment-free state has lost rows only to input events,
since its output actions use no commitments and retire no rows.  If
its current input is $i$, then the only output paired with $i$ is $f(i)$, and
every valid $j\succeq_I i$ has $f(i)\OutputOrd f(j)$, so $f(i)$ is exposable.
Conversely, suppose every output at every commitment-free reachable state is
exposable.  At $\sigma_0(i)$, exposability of $f(i)$ preserves the row
$(j,f(j))$ for each $j\succeq_I i$, so $f(i)\OutputOrd f(j)$ and $f$ is
monotone.
\end{proof}

Functional CALM and this proposition give the operational boundary.  A
monotone functional specification can emit every current output without a
commitment.  If monotonicity fails at $i\prec_I k$, the canonical output action
for $f(i)$ includes a commitment that removes the conflicting row at $k$.

Historically, CALM stands for \emph{Consistency As Logical Monotonicity}, and
its C has been read both as consistency and as confluence;
Definition~\ref{def:consistency} is both.  Functional CALM states the result
for any functional specification, regardless of the formalism used to express
it.  Before dropping the functional dependency, we give this consistency
condition a distributed interpretation.  Section~\ref{sec:distributed-execution}
realizes runs as participant runs; the construction will apply unchanged when
Section~\ref{sec:ambiguous-specifications} introduces output choice.
